\Askhsh{16804}{2}
\begin{ylh}{1}
    \item 9.1. Ορθές προβολές
    \item 9.4. Γενίκευση του Πυθαγόρειου θεωρήματος.
\end{ylh}
\begin{wrapfigure}[11]{r}{4cm} % Προσαρμόστε τον αριθμό γραμμών [11] εάν το κείμενο υπερβαίνει
\begin{tikzpicture}[scale=0.7] % Κλίμακα για καλύτερη προσαρμογή
    % Ορισμός σημείων με συντεταγμένες που να ταιριάζουν με τη γενική εμφάνιση της εικόνας
    % B στην αρχή, Γ στον άξονα x, A πάνω και αριστερά του H
    \tkzDefPoint(0,0){B}
    \tkzDefPoint(7,0){G} % Γ
    \tkzDefPoint(2,5){A} % A είναι το (2,5) για ένα οπτικό παράδειγμα, οδηγώντας στο H=(2,0)
    
    % Σχεδίαση του τριγώνου
    \tkzDrawPolygon(A,B,G)
    
    % Ύψος AH
    \tkzDefPointBy[projection=onto line B--G](A)
    \tkzGetPoint{H}
    \tkzDrawSegment[dashed](A,H)
    \tkzMarkRightAngle[size=0.2](A,H,G) % Σήμα ορθής γωνίας στο H
    
    % Ύψος BΘ
    \tkzDefPointBy[projection=onto line A--G](B)
    \tkzGetPoint{Th} % Θ
    \tkzDrawSegment[dashed](B,Th)
    \tkzMarkRightAngle[size=0.2](B,Th,G) % Σήμα ορθής γωνίας στο Θ
    
    % Σχεδίαση των σημείων
    \tkzDrawPoints(A,B,G,H,Th)
    
    % Ετικέτες
    \tkzLabelPoint[above](A){A}
    \tkzLabelPoint[below left](B){B}
    \tkzLabelPoint[below right](G){$\Gamma$}
    \tkzLabelPoint[below](H){H}
    \tkzLabelPoint[above right](Th){$\Theta$} % Τοποθετημένο για να ταιριάζει με την εικόνα
\end{tikzpicture}
\end{wrapfigure}
\Askhseis{ΘΕΜΑ 2}

Στο παρακάτω τρίγωνο $ΑΒ\Gamma$ φέρουμε τα ύψη του $ΑΗ$ και $Β\Theta$.

\begin{alist}
\item α) Να συμπληρώσετε τα κενά στις παρακάτω προτάσεις:
  \begin{rlist}
  \item Η προβολή της πλευράς $Β\Gamma$ στην πλευρά $Α\Gamma$ είναι το τμήμα $\Theta\Gamma$.
  \item Η προβολή της πλευράς $ΑΒ$ στην πλευρά $Β\Gamma$ είναι το τμήμα $BH$.
  \item Το τμήμα $Η\Gamma$ είναι η προβολή της πλευράς $A\Gamma$ στην πλευρά $B\Gamma$.
  \item Το τμήμα $Α\Theta$ είναι η προβολή της πλευράς $AB$ στην πλευρά $A\Gamma$.
  \item $Α\Gamma^2 = ΑΒ^2 + Β\Gamma^2 - 2 \cdot Β\Gamma \cdot BH$.
  \item $Β\Gamma^2 = AB^2 + Α\Gamma^2 - 2 \cdot A\Gamma \cdot Α\Theta$. \monades{15}
  \end{rlist}
\item β) Αν $ΑΒ = 4$, $Β\Gamma=5$ και $Α\Gamma = 6$, να υπολογίσετε το μήκος του τμήματος $Α\Theta$. \monades{10}
\end{alist}